\section{Herramientas utilizadas}
\label{sec:herramientas}

La tabla~\ref{tab:herramientas} indica qué herramienta resolvió cada parte del proyecto. Donde dice ``propio'', el algoritmo está escrito por el equipo sobre operaciones básicas de NumPy.

\begin{table}[ht]
\centering
\small
\begin{tabularx}{\linewidth}{@{}l X@{}}
\toprule
Herramienta & Qué resolvió \\
\midrule
Python 3.12 + NumPy \cite{harris2020} & Lenguaje del servidor y operaciones vectorizadas sobre imágenes. Base de todos los algoritmos propios. \\
FastAPI \cite{fastapi} + Uvicorn & Servidor HTTP y API; también sirve la interfaz web. \\
OpenCV \cite{bradski2000} & Lectura y escritura de imágenes, conversiones de color, desenfoques, morfología, \emph{distance transform}, clasificador de Haar, filtro guiado e \emph{inpainting} de Telea (solo como referencia de comparación). \\
MediaPipe \cite{lugaresi2019} & Modelos preentrenados, usados solo para inferencia: malla facial de 478 puntos, segmentador multiclase y puntos de la mano. También corre en el navegador para seguir la cabeza en la prueba en vivo de la peluca. \\
Ultralytics YOLOv8 + PyTorch \cite{jocher2023,paszke2019} & Ajuste fino y ejecución del segmentador de pelo propio (\archivo{hair_best.pt}). Es opcional: sin él, la aplicación usa MediaPipe. \\
Google Colab & Entrenamiento de YOLO con GPU (\archivo{backend/training/barberstudio_colab.ipynb}). \\
WebGL2 \cite{khronos2017} & Visor de gaussianas 3D en el navegador (\emph{shaders} propios). \\
HTML, CSS y JavaScript nativos & Interfaz sin \emph{frameworks} ni compilación. \\
pytest + GitHub Actions & 53 pruebas automáticas que se ejecutan en cada cambio. \\
Docker & Imagen reproducible con todas las dependencias y los modelos. \\
\LaTeX{} & Este documento. \\
\bottomrule
\end{tabularx}
\caption{Herramientas y la parte del proyecto que resolvió cada una.}
\label{tab:herramientas}
\end{table}

\begin{table}[ht]
\centering
\small
\begin{tabularx}{\linewidth}{@{}X l@{}}
\toprule
Pieza & Origen \\
\midrule
Detección de cara (Haar), 478 puntos faciales, segmentación base, puntos de la mano & Biblioteca o modelo \\
Segmentador de pelo YOLOv8 & Arquitectura de biblioteca, ajuste fino propio \\
Triangulación de Delaunay (Fortune y Bowyer--Watson) & Propio \\
\emph{Warp} afín, coordenadas baricéntricas, interpolación bilineal & Propio \\
Composición de la barba, hebras, \emph{feather}, anclaje a la malla & Propio \\
Tinte en CIELAB, tinte de barba, forma de la cara & Propio \\
Motor de cortes (\emph{pull-push}, cráneo, línea del pelo, sombreado) & Propio \\
Física de hebras (Verlet + PBD) y juego & Propio \\
Gaussianas 2D (gradientes analíticos, Adam) & Propio \\
Gaussianas 3D, proyección EWA, ordenamiento, exportación \texttt{.ply} & Propio \\
Peluca multivista (\emph{visual hull} por cortes, textura, estudio) & Propio \\
\bottomrule
\end{tabularx}
\caption{Qué viene de bibliotecas y qué implementó el equipo.}
\label{tab:propio}
\end{table}

Los efectos animados de la interfaz (texto que rota, títulos que aparecen, franja de líneas) son versiones en JavaScript nativo de componentes de React Bits, con licencia MIT y \emph{Commons Clause}. Los retratos de prueba de las figuras son retratos oficiales de la NASA, de dominio público.
